\documentclass[11pt]{letter}
\usepackage[T1]{fontenc}
\usepackage{newtxtext}
\usepackage[margin=25mm]{geometry}
\usepackage[colorlinks=true,urlcolor=blue]{hyperref}
%% DRAFT -- to be edited and signed by the corresponding author.
%% All blocks (date, closing, signature) are left-aligned.
\longindentation=0pt
\indentedwidth=\textwidth
\signature{Dr.\ Maohua Lin\\Corresponding author, on behalf of all authors\\Department of Biomedical Engineering, Florida Atlantic University\\777 Glades Road, Boca Raton, FL 33431, USA\\\href{mailto:mlin2014@fau.edu}{mlin2014@fau.edu}}
\begin{document}
\begin{letter}{The Editors\\Medical Image Analysis}
\opening{Dear Editors,}

We submit the manuscript ``GCZ-RL: Geometry-Calibrated Residual Learning for Spinal Ligament Landmark Localization on Vertebral Surface Meshes'' for consideration as a regular article in \emph{Medical Image Analysis}.

Subject-specific spine models need the attachment sites of the spinal ligaments, which are invisible on CT and are still placed by hand or by fixed geometric rules. The manuscript addresses the underlying image-analysis problem: localizing sparse ligament landmarks on vertebral surface meshes from a small set of expert annotations. The proposed method, GCZ-RL, learns ligament-specific admissible regions from annotations, transfers an atlas pattern by rigid registration, and learns only a bounded residual correction that is projected back onto the admissible region. Every output is therefore a surface vertex at an anatomically admissible site.

The evaluation was designed to be conservative. The primary analysis is a patient-level leave-one-subject-out protocol on every annotated lumbar vertebra, with all calibration and model selection nested inside the training subjects and a data-leakage test with corrupted held-out labels. Against the geometry-only baseline, GCZ-RL reduced the macro-averaged error from 2.879 to 2.150~mm, with improvement for every held-out subject and a 95\% bootstrap interval excluding zero. Against the released implementation of the dedicated SLD detector, run on the same vertebrae and scored with the same evaluator, accuracy was comparable (2.150 vs.\ 2.224~mm, not significant) with complementary class-specific strengths. We report coverage and matching-independent metrics, including SLD's own closest-point protocol, under which SLD is ahead, and we disclose every post hoc design decision.

All code, trained models, predicted landmarks, and the scripts that regenerate every table and figure are released, and the public source data are retrieved and verified by checksum. Parts of the method appear in the first author's doctoral dissertation; the manuscript has not been published or submitted elsewhere. All authors have approved the manuscript and agree with its submission. %% [AUTHORS: confirm, in particular Dr. Lin's agreement]

\closing{Best regards,}
\end{letter}
\end{document}
